\chapter{Conclusão}
\label{cap:conclusao}

\section{Síntese do trabalho}

O Cardly consolidou um ecossistema de estudo mobile integrado a API REST, atendendo ao objetivo central da disciplina: demonstrar domínio da integração React Native (Expo/TypeScript) com Spring Boot (Java) em arquitetura cliente-servidor \cite{react_native_docs,spring_boot_docs,tanenbaum2015}.

Principais entregas técnicas:
\begin{itemize}
    \item Autenticação JWT com suporte a Google SSO e perfil superadministrador \cite{jwt_io,google_oauth}.
    \item CRUD completo de assuntos e cartas com busca paginada via Specification \cite{jpa_specification}.
    \item Regras de revisão espaçada implementadas no \texttt{CardService}, com intervalos de 1 a 9 dias \cite{anki_manual,medium_spaced_repetition}.
    \item Front-end com NativeWind, flip animado, toasts e modais \cite{nativewind_docs,w3schools_flip}.
    \item Preparação para deploy em Railway e build mobile via EAS \cite{railway_docs,expo_eas_build}.
\end{itemize}

\section{Atendimento aos critérios de avaliação}

Conforme mapeamento do Capítulo~\ref{cap:objetivos} (Tabela~\ref{tab:criterios-avaliacao}), o projeto cobre:
\begin{itemize}
    \item Integração front-back: evidenciada nos capítulos de arquitetura, backend e frontend.
    \item Segurança: JWT, BCrypt, roles e SecureStore (Capítulo~\ref{cap:seguranca}).
    \item Organização de código: padrões 1.0, 1.1 e 2.0 documentados.
    \item UX: sessão de estudo fluida, paleta consistente, feedback ao usuário.
    \item Produção: variáveis de ambiente, migrations e checklist de deploy.
\end{itemize}

\section{Limitações identificadas}

\begin{enumerate}
    \item Calendário visual no dashboard ainda parcial (RF13).
    \item Módulo de amizades não implementado (RF14).
    \item Ausência de refresh token JWT; sessão expira em uma hora.
    \item Cobertura automatizada limitada no front-end mobile.
    \item Rate limiting e monitoramento avançado não implementados.
\end{enumerate}

Limitações são aceitáveis no recorte de MVP acadêmico, mas devem ser explicitadas na arguição.

\section{Trabalhos futuros}

\begin{itemize}
    \item Concluir calendário estilo Clue com heatmap mensal \cite{clue_app}.
    \item Implementar solicitações de amizade e compartilhamento social de decks.
    \item Adicionar notificações push (Expo Notifications) para cartas vencidas.
    \item Expandir testes E2E com Detox ou Maestro.
    \item Publicar OpenAPI via springdoc para documentação interativa \cite{openapi_springdoc}.
    \item Modo escuro completo utilizando tokens \texttt{background.dark}.
    \item Sincronização offline com fila de respostas (desafio avançado mobile).
\end{itemize}

\section{Aprendizados da equipe}

Estudantes do sétimo semestre reportaram curva de aprendizado significativa em:
\begin{itemize}
    \item Configuração de JWT end-to-end \cite{youtube_jwt_explained,stackoverflow_jwt_spring}.
    \item Animações nativas com Reanimated \cite{reanimated_docs,youtube_expo_typescript}.
    \item Specification JPA para filtros dinâmicos \cite{jpa_specification}.
    \item Disciplina de variáveis de ambiente e deploy \cite{twelve_factor,railway_docs}.
\end{itemize}

Vídeos e artigos citados ao longo deste relatório \cite{youtube_spring_boot_api,medium_spaced_repetition} foram recursos pedagógicos essenciais.

\section{Considerações finais}

O Cardly demonstra que flashcards com revisão espaçada constituem domínio adequado para aplicação mobile com regras de negócio não triviais, indo além de CRUD superficial exigido em projetos introdutórios. A documentação formal em ABNT \cite{abnt_nbr14724,abnt_nbr6023}, aliada à demonstração funcional integrada, prepara a equipe para apresentação prática e arguição técnica previstas na metodologia da disciplina.

A versão final deste relatório depende da inserção de figuras reais (prints, diagramas), revisão ortográfica por três revisores (dois com modelo Opus conforme planejamento interno), dados do orientador na capa e compilação PDF via Overleaf ou LaTeX local com fonte Times New Roman conforme manual UNILESTE.

\section{Referências aos repositórios}

Código-fonte e histórico de desenvolvimento:
\begin{itemize}
    \item Backend: \url{https://github.com/hugoFreit4s/cardly-backend} \cite{cardly_backend_repo}
    \item Frontend: \url{https://github.com/hugoFreit4s/cardly-rnative-app} \cite{cardly_frontend_repo}
\end{itemize}
